\par\medskip
\hypertarget{native-086-article067}{}\pdfbookmark[2]{67.}{article-067}
\begin{center}67.\end{center}
Methodus alia, problema art. praec. solvendi, ex aequationibus\HistoricalNoteEight{Handschriftliche Bemerkung:]
\[
\begin{aligned}
\sin\alpha\cos\delta&=-\sin\epsilon\sin b+\cos\epsilon\cos b\sin l\\[.35em]
\cos\alpha\cos\delta&=\cos b\cos l\\[.35em]
\sin\delta&=\cos\epsilon\sin b+\sin\epsilon\cos b\sin l.
\end{aligned}
\]
}
\[
\begin{aligned}
\cos\epsilon\sin l&=\sin\epsilon\operatorname{tang}b+\cos l\operatorname{tang}\alpha\\[.4em]
\sin\delta&=\cos\epsilon\sin b+\sin\epsilon\cos b\sin l\\[.4em]
\cos b\cos l&=\cos\alpha\cos\delta
\end{aligned}
\]
petitur. Determinetur angulus auxiliaris $\theta$ per aequationem
\[
\operatorname{tang}\theta=\frac{\operatorname{tang}b}{\sin l},
\]
eritque
\[
\begin{aligned}
\operatorname{tang}\alpha&=\frac{\cos(\epsilon+\theta)\operatorname{tang}l}{\cos\theta}\\[.6em]
\operatorname{tang}\delta&=\sin\alpha\operatorname{tang}(\epsilon+\theta),
\end{aligned}
\]
\SourceLeaf{087}quibus aequationibus ad calculi confirmationem adiici potest
\[
\cos\delta=\frac{\cos b\cos l}{\cos\alpha}\quad\text{sive}\quad
\cos\delta=\frac{\cos(\epsilon+\theta)\cos b\sin l}{\cos\theta\sin\alpha}\cdot
\]
Ambiguitas in determinatione ipsius $\alpha$ per aequ. secundam eo tollitur, quod $\cos\alpha$ et $\cos l$ eadem signa habere debent.

Haec methodus minus expedita est, si praeter $\alpha$ et $\delta$ etiam $E$ desideratur: formula commodissima ad hunc angulum determinandum tunc erit $\cos E=\frac{\sin\epsilon\cos\alpha}{\cos b}=\frac{\sin\epsilon\cos l}{\cos\delta}\cdot$ Sed per hanc formulam $E$ accurate computari nequit, quoties $\pm\cos E$ parum ab unitate differt; praeterea ambiguitas remanet, utrum $E$ inter $0$ et $180^0$ an inter $180^0$ et $360^0$ accipere oporteat. Incommodum prius raro ullius momenti est, praesertim, quum ad computandas rationes differentiales ultima praecisio in valore ipsius $E$ non requiratur: ambiguitas vero illa adiumento aequationis $\cos b\cos\delta\sin E=\cos\epsilon-\sin b\sin\delta$ facile tollitur, quae ostendit $E$ inter $0$ et $180^0$ vel inter $180^0$ et $360^0$ accipi debere, prout $\cos\epsilon$ maior fuerit vel minor quam $\sin b\sin\delta$: manifesto hoc examen ne necessarium quidem est, quoties alteruter angulorum $b$, $\delta$ limitem $66^0\,32'$ non egreditur: tunc enim $\sin E$ semper fiet positivus. Ceterum eadem aequatio in casu supra addigitato ad determinationem exactiorem ipsius $E$, si operae pretium videtur, adhiberi poterit.
